\documentclass[11pt, a4paper, oneside]{article}
\usepackage[ngerman]{babel}
\usepackage{worksheet}

\begin{document}
	\author{L. Bung}
	\title{Exkurs: Satz von Vieta}
	\subject{Mathematik}
	\class{FHR1}
	\maketitle
	
	\begin{hint}{Der Satz von Vieta}
		Der \textbf{Satz von Vieta} (benannt nach Fran\c{c}ois Vieta, 1540--1603) ist eine Möglichkeit, um die Nullstellen von gemischt-quadratischen Gleichungen schnell herauszufinden:
		
		\begin{quote}
			Sind $x_1$ und $x_2$ Nullstellen der quadratischen Funktion $f(x) = a(x^2+px+q)$, so gilt $p = -(x_1 + x_2)$ und $q = x_1 \cdot x_2$.
		\end{quote}
		
		Die Funktion $f(x) = 2x^2-10x+12$ kann auch geschrieben werden als $f(x) = 2(x^2-5+6)$.
		Man sieht: $p = -(2+3)$ und $q = 2 \cdot 3$.
		Also sind die Nullstellen der Funktion $x_1 = 2$ und $x_2 = 3$.
	\end{hint}
	
	\singletask{Nullstellen berechnen}
	
	Nutzen Sie den Satz von Vieta, um die Nullstellen der folgenden Funktionen zu bestimmen:
	
	\begin{multicols}{3}
		\begin{enumerate}[label=\alph*)]
			\item $a(x) = 3x^2+6x-9$
			\item $b(x) = -\frac{1}{2}x^2-3x-4$
			\item $c(x) = \frac{2}{3}x^2+\frac{2}{3}x-\frac{4}{3}$
		\end{enumerate}
	\end{multicols}
	
	\checkered[11cm]
	
	\let\clearpage\relax
\end{document}